\documentclass[11pt]{amsart}
\usepackage[T1]{fontenc}
\usepackage{lmodern}
\usepackage{amsmath,amssymb,amsthm}
\usepackage{xcolor}
\usepackage[colorlinks=true,linkcolor=blue!55!black,citecolor=blue!55!black,urlcolor=blue!55!black]{hyperref}
\newtheorem{theorem}{Theorem}
\newtheorem{proposition}[theorem]{Proposition}
\title[A 201-Point Folkman Finite-Sums Colouring]{A 201-Point Two-Colouring for the Four-Generator Folkman Finite-Sums Problem}
\author{Zijian Zeng}
\address{Institute of Computer Science and Digital Innovation, UCSI University, Kuala Lumpur, 56000, MALAYSIA}
\email{zijianzeng@foxmail.com}
\email{1002266693@ucsiuniversity.edu.my}
\date{September 1, 2026}
\hypersetup{pdftitle={A 201-Point Two-Colouring for the Four-Generator Folkman Finite-Sums Problem},pdfauthor={Zijian Zeng}}

\begin{document}
\begin{abstract}
Let $F(k)$ be the least $N$ such that every two-colouring of $[N]$ has a
$k$-element set whose nonempty subset sums lie in $[N]$ and are all one
colour.  We give an explicit two-colouring of $[201]$ with no such
four-element set.  Two independent exhaustive verifiers check all
$2{,}587{,}336$ eligible generator sets using exact integer arithmetic.  Thus
$F(4)\geq202$, disproving the proposed equality $F(4)=201$.  We do not
determine $F(4)$.
\end{abstract}
\maketitle

\section{Introduction}
For a finite set $A$ of positive integers, put
\[
 \operatorname{FS}(A)=\left\{\sum_{a\in B}a:\varnothing\ne B\subseteq A\right\}.
\]
Folkman's theorem ensures that the parameter in the abstract is finite.  The
standard quantitative formulation appears explicitly in
\cite{BaloghEtAl2017}; later work records the lower bound of that paper and
Taylor's tower-type upper bound as the then-known asymptotic frontier
\cite[p.~5]{BaloghLavrovShakanWagner2019}.

A public computational report dated July 26, 2026 gave an explicit colouring
of $[200]$. It checked $2{,}535{,}064$ eligible generator sets and thereby
proved $F(4)\geq201$ \cite{WorkingReport531}.  We extend that concrete
finite record by one point.  Searches recorded with the accompanying data
found no earlier primary or official record of the bound below as of
September 1, 2026; this is a bounded literature statement, not a claim about
uncatalogued work.

\section{The finite reduction}
Let $H(N,4)$ be the hypergraph with vertex set $[N]$ and edges
\[
 \operatorname{FS}(A),\qquad A\in\binom{[N]}4,\quad \sum_{a\in A}a\leq N.
\]
The sum condition is equivalent to $\operatorname{FS}(A)\subseteq[N]$ because
the generators are positive.  Subset-sum collisions are deduplicated: edges
are sets and may have fewer than $15$ vertices.  Directly from the definition,
$F(4)$ is the first $N$ for which $H(N,4)$ is not two-colourable.

\section{The colouring}
Concatenate the following five lines and read the resulting bits as the colours
of $1,\ldots,201$:
\begin{verbatim}
01110111101101100001110011000000100100000100011100
10011010001011110101000110010100000101111101101010
10011110001111110100100110111010101011111110101001
10100010100100101111110000100011100011001010000011
1
\end{verbatim}

\begin{theorem}\label{thm:main}
The displayed word properly two-colours $H(201,4)$.  Consequently
$F(4)\geq202$; in particular, $F(4)\ne201$.
\end{theorem}

\begin{proof}
The exact finite check is summarized in Proposition~\ref{prop:verification}.
Every eligible edge meets both colour classes, which is precisely proper
two-colourability of $H(201,4)$.  Hence the universal forcing property fails
at $N=201$, so the integer $F(4)$ is at least $202$.
\end{proof}

The construction is especially transparent relative to the $200$-point
prefix.  There are $52{,}272$ new generator sets with total exactly $201$.
Under the old prefix, only $\{38,41,43,79\}$ has all its other subset sums in
one colour, namely colour $0$; no new set gives the opposite requirement.
Thus colour $1$ at $201$ extends the prefix.  The theorem nevertheless rests
on the full verification below, not on trusting the prefix.

\section{Independent exhaustive verification}
\begin{proposition}\label{prop:verification}
Two independent definition-level enumerations both check all $2{,}587{,}336$
eligible sets and find no monochromatic finite-sums edge.  Their common
edge-size histogram is
\[
\begin{array}{c|rrrrrr}
|\operatorname{FS}(A)|&10&11&12&13&14&15\\ \hline
\#A&20&1585&2952&151189&61964&2369626.
\end{array}
\]
\end{proposition}
\begin{proof}
The first verifier iterates over all four-combinations of $[201]$ and obtains
each edge by enumerating every nonempty subcombination.  The second uses four
explicit strictly ordered loops with independent upper bounds and obtains the
edge by iterative sumset expansion, starting from $\{0\}$ and adjoining one
generator at a time.  Each program parses exactly $201$ binary digits, rejects
malformed input, constructs mathematical sets to handle collisions, and stops
on the first monochromatic edge.  Both terminate successfully with the counts
shown above.  These exhaustive integer computations constitute a directly
checkable finite certificate.
\end{proof}

\section{Reproducibility and limitations}
From the source archive, run
\begin{verbatim}
python3 code/verify_witness_a.py
python3 code/verify_witness_b.py
\end{verbatim}
with Python 3.10 or later.  The certificate is stored as
\url{certificate/f4_n201_coloring.txt}; neither verifier requires a SAT
solver or a nonstandard Python package.  Discovery and verification design
are separated from the theorem check.  The mathematical claim is bound to the
displayed exact data and independently executable source; no uncheckable
solver verdict is used.

This result settles only the stated finite cell and rules out one proposed
exact value.  It gives no upper bound at $202$ and does not determine $F(4)$.

\bibliographystyle{amsplain}
\bibliography{references}
\end{document}
